\SourceLeaf{109}%
nen Resultaten eine grössere Genauigkeit beilegt, als sie wirklich besitzen. Freilich hat in dem Falle, wo die Anzahl der Beobachtungen vielemale grösser ist als die der unbekannten Grössen, diese Unrichtigkeit wenig zu bedeuten; allein theils erfordert die Würde der Wissenschaft, dass man vollständig und bestimmt übersehe, wieviel man hierdurch zu fehlen Gefahr läuft, theils sind auch wirklich öfters nach jenem fehlerhaften Verfahren Rechnungsresultate in wichtigen Fällen aufgestellt, wo jene Voraussetzung nicht Statt fand. Der Verfasser hat daher diesen Gegenstand einer besondern Untersuchung unterworfen, die zu einem sehr merkwürdigen höchst einfachen Resultate geführt hat. Man braucht nemlich den nach dem angezeigten fehlerhaften Verfahren gefundenen mittlern Fehler, um ihn in den richtigen zu verwandeln, nur mit
\[
\surd\frac{\pi-\rho}{\pi}
\]
zu multipliciren, wo \(\pi\) die Anzahl der Beobachtungen und \(\rho\) die Anzahl der unbekannten Grössen bedeutet.

Die letzte Untersuchung betrifft noch die Ausmittelung des Grades von Genauigkeit, welcher dieser Bestimmung des mittlern Fehlers selbst beigelegt werden muss: die Resultate derselben müssen aber in der Abhandlung selbst nachgelesen werden.

\par\bigskip
\begin{center}
\rule{.88\textwidth}{.3pt}\par
Göttingische gelehrte Anzeigen.\quad 1826 September 25.\par
\rule{.88\textwidth}{.3pt}
\end{center}
\pdfbookmark[1]{1826 September 25}{anzeige-1826-09-25}

Am 16. September überreichte der Herr Hofr. \textsc{Gauss} der königl. Societät eine Vorlesung:
\begin{center}
\textit{Supplementum Theoriae combinationis observationum erroribus minimis obnoxiae.}
\end{center}
\noindent Bei allen frühern Arbeiten über die Anwendung der Wahrscheinlichkeitsrechnung auf die zweckmässigste Benutzung der Beobachtungen, und namentlich auch in der Behandlung dieses Gegenstandes im fünften Bande der \textit{Commentationes recentiores}, liegt in Beziehung auf die Form der Hauptaufgabe eine bestimmte Voraussetzung zum Grunde, die allerdings den meisten in der Ausübung vorkommenden Fällen angemessen ist. Diese Voraussetzung besteht darin, dass die be\SourceLeaf{110}obachteten Grössen auf eine bekannte Art von gewissen unbekannten Grössen (Elementen) abhängen, d. i. bekannte Functionen dieser Elemente sind. Die Anzahl dieser Elemente muss, damit die Aufgabe überhaupt hierher gehöre, kleiner sein, als die Anzahl der beobachteten Grössen, also diese selbst abhängig von einander.

Inzwischen sind doch auch die Fälle nicht selten, wo die gedachte Voraussetzung nicht unmittelbar Statt findet, d. i. wo die beobachteten Grössen noch nicht in der Form von bekannten Functionen gewisser unbekannter Elemente gegeben sind, und wo man auch nicht sogleich sieht, wie jene sich in eine solche Form bringen lassen; wo hingegen zum Ersatz die gegenseitige Abhängigkeit der beobachteten Grössen (die natürlich auf irgend eine Weise gegeben sein muss) durch gewisse Bedingungsgleichungen gegeben ist, welchen die wahren Werthe von jenen, der Natur der Sache nach, nothwendig genau Genüge leisten müssen. Zwar sieht man bei näherer Betrachtung bald ein, dass dieser Fall von dem andern nicht wesentlich, sondern bloss in der Form verschieden ist, und sich wirklich, der Theorie nach leicht, auf denselben zurückführen lässt: allein häufig bleibt dies doch ein unnatürlicher Umweg, der in der Anwendung viel beschwerlichere Rechnungen herbeiführt, als eine eigne der ursprünglichen Gestalt der Aufgabe besonders angemessene Auflösung. Diese ist daher der Gegenstand der gegenwärtigen Abhandlung, und die Auflösung der Aufgabe, welche sie als ein selbstständiges von der frühern Abhandlung unabhängiges Ganze gibt, hat ihrerseits eine solche Geschmeidigkeit, dass es sogar in manchen Fällen vortheilhaft sein kann, sie selbst da anzuwenden, wo die bei der ältern Methode zum Grunde liegende Voraussetzung schon von selbst erfüllt war.

Die Hauptaufgabe stellt sich hier nun unter folgender Gestalt dar. Wenn von den Grössen \(v,\,v',\,v''\) u. s. w., zwischen welchen ein durch eine oder mehrere Bedingungsgleichungen gegebener Zusammenhang Statt findet, eine andere auf irgend eine Art abhängig ist, z. B. durch die Function \(u\) ausgedrückt werden kann, so wird eben dieselbe auch auf unendlich viele andere Arten aus jener bestimmt, oder durch unendlich viele andere Functionen, statt \(u\), ausgedrückt werden können, die aber natürlich alle einerlei Resultate geben, in so fern die wahren Werthe von \(v,\,v',\,v''\) u. s. w., welche allen Bedingungsgleichungen Genüge leisten, substituirt werden. Hat man aber nur genäherte Werthe von \(v,\,v',\,v''\) u. s. w., wie sie Beobachtungen von beschränkter Genauigkeit immer nur liefern können, so
\SourceLeaf{111}
können auch die daraus abgeleiteten Grössen auf keine absolute Richtigkeit Anspruch machen: die verschiedenen für \(u\) angewandten Functionen werden, allgemein zu reden, ungleiche, aber was die Hauptsache ist, ungleich zuverlässige Resultate geben. Die Aufgabe ist nun, aus der unendlichen Mannigfaltigkeit von Functionen, durch welche die unbekannte Grösse ausgedrückt werden kann, diejenige auszuwählen, bei deren Resultat die möglich kleinste Unzuverlässigkeit zu befürchten bleibt.

Die Abhandlung gibt eigentlich \textit{zwei} Auflösungen dieser Aufgabe. Die erste Auflösung erreicht das Ziel auf dem kürzesten Wege, wenn wirklich nur Eine unbekannte von den Beobachtungen auf eine vorgeschriebene Art abhängige Grösse abzuleiten ist. Allein die nähere Betrachtung dieser Auflösung führt zugleich auf das merkwürdige Theorem, dass man für die unbekannte Grösse genau denselben Werth, welcher aus der zweckmässigsten Combination der Beobachtungen folgt, erhält, wenn man an die Beobachtungen gewisse nach bestimmten Regeln berechnete Veränderungen anbringt, und sie dann in irgend einer beliebigen Function, welche die unbekannte Grösse ausdrückt. substituirt. Diese Veränderungen haben neben der Eigenschaft, dass sie allen Bedingungsgleichungen Genüge leisten, noch die, dass unter allen denkbaren Systemen, welche dasselbe thun, die Summe ihrer Quadrate (in so fern die Beobachtungen als gleich zuverlässig vorausgesetzt wurden) die möglich kleinste ist. Man sieht also, dass hierdurch zugleich eine neue Begründung der Methode der kleinsten Quadrate gewonnen wird, und dass diese von der Function \(u\) ganz unabhängige \textit{Ausgleichung} der Beobachtungen eine zweite Auflösungsart abgibt, die vor der ersten einen grossen Vorzug hat, wenn mehr als Eine unbekannte Grösse aus den Beobachtungen auf die zweckmässigste Art abzuleiten ist: in der That werden die Beobachtungen dadurch zu \textit{jeder} von ihnen zu machenden Anwendung fertig vorbereitet. Nur musste bei dieser zweiten Auflösung noch eine besondere Anleitung hinzukommen, den Grad der Genauigkeit, der bei jeder einzelnen Anwendung erreicht wird, zu bestimmen. Für dies alles enthält die Abhandlung vollständige und nach Möglichkeit einfache Vorschriften. die natürlich hier. keines Auszuges fähig sind. Eben so wenig können wir hier in Beziehung auf die, nach der Entwickelung der Hauptaufgaben, noch ausgeführten anderweitigen Untersuchungen, welche mit dem Gegenstande in innigem Zusammenhange stehen, uns in das Einzelne einlassen. Nur das Eine merkwürdige Theorem führen wir hier an, dass
